\documentclass[10pt]{article}
\usepackage{lingmacros}
\usepackage{tree-dvips}
\begin{document}
\pagenumbering{gobble}
\section*{1. Classificazione microrganismi}
\textbf{Saprofiti} come miceti che decompongono materiale organico rendendolo disponibile alle piante.\\
\textbf{Simbionti} come microrganismi che risiedono nel rumine dei ruminanti, gli azotofissatori delle leguminose e i licheni sulle alghe fotosintetiche/cianobatterio.\\
\textbf{Commensali} come microbiota intestinale e possono essere opportunisti (candida).\\
\textbf{Parassiti} come Colera (vibrione flagellato).\\
\textbf{Elminti} come vermi e tenie. \\
\section*{2. Malattie infettive}
L'\textbf{infezione} avviene principalmente sulle mucose del corpo (naso, bocca, ano, occhi).\\
L'\textbf{infestazione} è l'invasione in un organismo di parassiti e/o animali (pidocchi, elminti, scabbia).\\
\textbf{Dose Infettante 50\%} (DI50) è la quantità di microrganismi infettanti necessari a infettare il 50\% delle cavie. \\
\textbf{DL50} è la quantità di microrganismi infettanti necessari a uccidere il 50\% delle cavie. \\
\paragraph*{Malattia di Lyme} Causata da una spirocheta portata da una \underline{zecca} che succhia il sangue dai cervi o roditori, solamente le femmine sono ematofaghe perché aiuta con la produzione di uova. Provoca eritemi e può aggravare con il tempo con artrite e carditi fino a demielinizzazione neuronale con AD o sclerosi multipla. Grazie al loro pungiglione le zecche usano un movimento rotatorio quindi devono essere rimosse con un movimento contrario, inoltre si deve evitare i disinfettanti perché causa il rigurgito della zecca. \\
\paragraph*{Legionellosi} Causata dalla Legionella pneumophila (gram- opportunista) ha un'incidenza maggiore negli anziani e si riproduce negli \underline{ambienti acquatici} che risale attraverso tubature e aria condizionata dove forma \underline{biofilm}, non ha trasmissione persona per persona. La Legionella previene il killing intracellulare lidendo il fagolisosoma dopo essersi riprodotta all'interno. \\\\
Un trattamento di antibiotici o chemio/radioterapici può causare una disbiosi intestinale portando a una rottura dell'equilibrio dell'ecosistema e favorendo lo sviluppo di candidosi. \\
Le peritoniti possono essere causate da un varco nell'intestino dovuto ad appendiciti che permette ai microrganismi intestinali a infiammare il peritoneo.\\
Le \textbf{zoonosi} possono fare un salto di specie diretto oppure usare un \underline{ospite intermedio} che entrando in contatto con un virus umano e di un'altra specie animale genera un virus ibrido. L'HIV è una zoonosi (spill-over) e nasce come virus delle scimmie.\\
La trasmissione di tubercolosi è dovuta all'espettorato e non semplice saliva, mentre altri microrganismi vengono trasmessi con \textbf{microgocce di Flugge}, le quali raggiungono dimensioni di 1-4 micron che gli permette di viaggiare in aria per molto tempo con al loro interno anche più di un solo microrganismo.\\
Il Treponema pallidum (sifilide) e la Neisseria gonorrhoeae (gonorrea) sono esclusivamente a trasmissione sessuale perché non sopravvivono facilmente all'esterno e usano il contatto tra mucose per trasmettersi. \\
L'infezione subclinica è un'infezione asintomatica come portatore sano. \\
Cellule M??
\subsubsection*{Microbiota} 
Arriva ad un peso complessivo di circa 2kg con un rapporto 10:1 rispetto alle cellule. Si trovano principalmente in cute, mucose e apparato digerente, in quest'ultimo prolifera nella bocca e nel colon piuttosto che nello stomaco dove si trovano solo lattobacilli e elicobacter pilori. Nella cute la distribuzione è concentrata nelle zone più umide.\\
Sono molto utili perché sintetizzano vitamina K e B, competono con microrganismi esterni e stimolano il sistema immunitario.\\
La composizione del microbiota è influenzato da età (sterili alla nascita e dopo un anno si stabilizza), dieta, ambiente e antibiotici (causando disbiosi). \\
Sulla cute troviamo lo stafilococco epidermidis (biofilm), nella bocca lo streptococcus mutans (carie). 
\paragraph*{Interferenza batterica} Il microbiota occupa i recettori, concorre nell'uso di sostanze nutritive, produce batteriocine (collicine dai E.Coli) e abbassano pH.\\\\
Il \textbf{lactobacillus acidophilus} prolifera nell'ambiente vaginale producendo acido lattico dal glicogeno abbassando il pH e proteggendo la vagina.  \paragraph*{Disbiosi} Analisi su \underline{topi germ-free} hanno evidenziato che l'assenza di microbiota porta una diminuzione di cellule staminali nell'intestino, vasi nei villi e del sistema immunitario, inoltre soffrono di stati di ansia e capacità relazionale ridotta (analizzato con dark light box), questo comportamento alterato si può ripristinare trapiantando microbiota fecale di una popolazione normale nei germ-free solo giovani. I batteri gram- bacterio detes e i gram+ firmicutes sono in equilibrio per il fisiologico, mentre nelle persone obese aumenta per il 20\% il numero di fermicutes che si occupano di liberare energia dagli alimenti. \\\\
Alcuni microrganismi producono acidi grassi a corta catena che hanno effetto anoressizzante nel cervello diminuendo la fame. Mentre l'iniezione di Campylobacter provoca stati d'ansia nei topi.
\paragraph*{Probiotici} Nei primi del '900 si è notata una vita longeva in coloro che avevano una dieta a prevalenza di yogurt e altri latticini che aiuta la proliferazione batterica e abbassa il pH impedendo lo sviluppo di batteri proteolitici (che producono fenoli e ammoniaca); da questo sono stati definiti i \textbf{probiotici} come microrganismi vivi che se assunti in \underline{adeguate quantità} garantiscono benefici. I probiotici devono essere di origine umana e devono resistere ad acidità gastrica e bile per raggiungere la mucosa intestinale, tra questi lactobacillus, bifidobacterio, streptococcus e saccharomyces boulardii (l'unico lievito). \\
Non devono essere presi in quantità eccessive e soprattutto non in condizioni di salute perché possono causare un disequilibrio intestinale.\\
Similmente i \textbf{prebiotici} sono sostanze alimentari che aumentano l'attività e la crescita di specifici microrganismi. \\\\
Il \textbf{trapianto di feci} può essere usato per patologie acute agendo con molta efficacia ma per breve periodo rispetto all'uso di antibiotici, un attecchimento più elevato avviene nel caso il trapianto provenga da consanguinei che possiedono mucose simili.
\section*{3. Difese}
\paragraph*{Psiconeuroendocrinoimmunologia} Lo stato psicologico dell'individuo influenza il sistema immunitario grazie all'interconnessione con il sistema endocrino.
\subsection*{Fagociti professionali}
La \textbf{fagocitosi} avviene tramite prolungamenti della membrana, gli \underline{pseudopodi}, che permettono di portare all'interno l'estraneo e formare il \underline{fagosoma}, quest'ultimo si fonde al lisosoma formando \underline{fagolisosoma} dove viene digerito. Nel caso di tubercolosi e legionella il batterio inibisce il killing nei fagosomi e il soggetto è infettato latentemente finché non si abbassano le difese, momento in cui il batterio si replica e causa malattia.\\
Il \textbf{killing} avviene con meccanismi Ossigeno indipendenti (lisozima e protelitici), se il microrganismo sopravvive si usano quelli Ossigeno dipendenti che usano NADPH ossidasi per la produzione di \underline{anione superossido} che è altamente citocida, chi è affetto da malattia granulomatosa cronica non presenta NADPH ossidasi. Alcuni batteri possiedono il \underline{superossido dismutasi} che produce ossigeno e acqua ossigenata, per eliminare l'acqua ossigenata usano catalasi. \\
Per resistere alla fagocitosi può essere prodotta una capsula, si può uccidere il fagocita e resistere al fagosoma crescendo al suo interno.\\ 
I \textbf{polimorfonucleati} o neutrofili hanno un nucleo polilobato e sono i leucociti dominanti nel sangue, prendono energia dalla \underline{glicolisi anaerobia} agendo in tessuti profondi dove c'è ipossia, sono ricchi di lisosomi e vivono circa 24h senza dividersi. Si occupano di microrganismi \underline{extracellulari piogeni (produttori di pus)}.\\
I \textbf{macrofagi} originano dai \underline{monociti} che circolano e raggiungono endoteli o connettivi dove si differenziano, hanno una vita di 4/8gg se attivi o anche anni da inattivi e agiscono in condizione di aerobiosi. Si trovano soprattutto in fegato (Kupfer), polmone (elimina corpuscolati del fumo), encefalo (microglia), milza e linfonodi, rene. Colpiscono microrganismi \underline{intracellulari} come salmonelle e legionelle.\\\\
La \textbf{difesa immediata} anatomica (cute e mucose) o chimiche (NK, fagociti e lisozima) non hanno memoria immunologica e si attivano in massimo 4h.\\
La \textbf{difesa innata} si basa sulla presenza di \underline{PAMP}, profili molecolari associati al patogeno, che vengono riconosciuti dai recettori (PRR) di PMN, macrofagi e cellule dendritiche, è abbastanza veloce ma non ha memoria immunologica. I PRR o Toll-like receptors (TLR) sono proteine transmembrana che all'interno liberano un fattore di trascrizione (NF-kB) quando riconoscono PAMP o anche \underline{DAMP}, segnali di stress endogeno come acido urico e ATP rilasciate per un danno ai tessuti; la cellula produce citochine (attivano infiammazione) e molecole costimolatorie nelle cellule fagocitiche presentanti antigene. Ogni \underline{TLR} (13 in totale) ha una specificità abbastanza ridotta per strutture diversa: il TLR-4 agisce sui gram- (lipopolisaccaride), TLR-3 su RNA a doppia catena (virus), TLR-5 sui flagelli. \\
La \textbf{difesa adattiva} è molto più lenta ma attiva i linfociti B e T per la produzione di memoria immunologica.\\
Le cellule \textbf{natural killer} agiscono sui tumori avvicinandosi alla cellula e liberando perforine e granzimi che creano pori sulla membrana e inducono l'apoptosi. Possiedono sulla membrana recettori che legano l'\underline{MHC 1} (istocompatibilità) nelle cellule normali, se non è presente si attiva, i virus possono deprimere l'espressione di MHC 1.
\paragraph*{Complemento} Componenti proteici (da C1 a C9) del siero labili al calore che formano il MAC, complesso di attacco alla membrana formando pori nell'estraneo. Le componenti sono prodotti da fegato, cervello, polmone. \\
L'attivazione per \textbf{via classica} (scoperta prima) avviene dopo la formazione dell'\underline{immunocomplesso} che attiva il C1, questo scinde C4 e C2 in C4a-b e C2a-b, il C4b e il C2a si associano a formare C3 convertasi che scinde il C3 per far associare il C3b all'enzima, questo forma la C5 convertasi producendo C5b che si deposita sul microrganismo che attira C6,7,8 e molti C9 che portano alla lisi dell'estraneo.\\
Nella \textbf{via alternativa} non è necessario l'immunocomplesso rientrando quindi nelle \underline{difese innate}. Il C3b lega al Fattore B, quest'ultimo viene scisso dal Fattore D producendo C3b-Bb convertasi, un \underline{complesso instabile} che si degrada se non sono presenti estranei, si stabilizza e attiva con componenti microbici riprendendo il percorso della via classica dalla formazione del C3 convertasi.\\
Tutti i componenti scissi non usati si solubilizzano nel siero come molecole \textbf{chemiotattiche} che guidano i PMN con chemiotassi sui microrganismi estranei.
\subsection*{Difesa immunitaria} 
La risposta immunitaria è \underline{specifica} per i diversi antigeni presenti, sono così tanti i linfociti presenti \underline{diversi} tra loro ($10^{14}$) che presentano recettori per antigeni presenti su microrganismi mai incontrati dall'individuo. La \underline{memoria} è fondamentale perché rende più efficiente la risposta di un'infezione già avvenuta, inoltre si \underline{autolimita} mettendo fine alla risposta dopo un po' di tempo ma non sempre avviene correttamente. Infine è in grado di discernere il self dal non self grazie alla \underline{tolleranza immunitaria}, cosa che non avviene nelle autoimmuni.
\paragraph*{Anticorpi}Difese umorali perché vengono secrete nel siero e mucose, sono proteine globulari, da cui immunoglobuline (Ig). I monomeri sono formati da due catene pesanti (H) e due leggere (L) tenute insieme da ponti disolfuro a formare una Y, le braccia (H+L) sono il paratopo, sito variabile, mentre l'estremità opposta (2H, frammento cristallizabile) funge da sito di legame con il fagocita permettendo di fare da ponte tra fagocita e antigene. I monomeri sono tenuti insieme dal peptide J.
\subparagraph*{IgG} Monomerica (75\% delle Ig sieriche), deriva da cellule della memoria, \underline{attravera la placenta} e lega con più forza. \underline{Neutralizza tossine} batteriche e virus impedendo di legarsi ad altre cellule.
\subparagraph*{IgM} Pentamerica, \underline{la prima prodotta} risposta primaria e quindi l'unica che viene prodotta dal \underline{neonato} nei primi giorni. Il linfocita B della memoria produce prima IgM e poi attraverso \underline{switching isotipico} produce IgG.
\subparagraph*{IgA} Monomerica (siero) e dimerica (muco), vengono eliminate con il \underline{muco} quindi sono in piccola quantità nel siero. Sono presenti nel \underline{colostro} formando l'immunità passiva naturale con IgG.
\subparagraph*{IgE} Monomerica, lega l'estremità 2H ai mastociti polmonari e intestinali stimolando il rilascio di \underline{istamina} per una risposta infiammatoria localizzata.
\subparagraph*{IgD} Monomerica, si trova sui linfociti B con le IgM \underline{regolandone l'attività}.
\\\\
\textbf{Opsonizzano} gli antigeni capsulati grazie a vaccini prodotti con gli zuccheri della capsula, permettendo l'attività di PMN e macrofagi. Attivano la via classica del complemento, fondamentale perché lo stesso C3b funziona come opsonina come le Ig.\\
La \textbf{sieroterapia} è un'immunità passiva artificiale l'inoculazione di Ig animali (cavallo, può causare allergie) o umane (vaccinati recentemente) specifiche per malattie come difterite, tetano, botulino, veleni e gangrena gassosa (Perfringens da proiettili). L'inoculazione di siero e vaccino devono essere fatti a distanza di tempo oppure il siero inculato colpisce il vaccino.
\section*{4. Cellula batterica}
Il mordenzante (Lugol) stabilizza il legame tra cellule e colorante.\\
Svedberg: coefficiente di sedimentazione di una particella che dipende dal peso.\\
Nella membrana plasmatica batterica non sono presenti steroli e proteine glicosilate, i carboidrati si trovano solo in glicolipi e glicosfingolipidi. È anche la zona dove avviene la fosforilazzione ossidativa.
\paragraph*{Strato LPS} (lipopolisaccaride) costituito da 3 strati, il lipide A è la porzione tossica interna, successivamente il polisaccaride del core e infine l'antigene O (somatico) che è diverso per ogni microrganismo, costituito da 4 zuccheri ripetuti più di 10 volte. L'LPS è un'endotossina che si libera durante la duplicazione e lisi dei gram-, se la quantità è eccessiva si genera uno \underline{shock endotossico}.\\
I macrofagi hanno un recettore CD14 non transmembrana specifico per il complesso LPS-LBP (LPS binding protein), per attivare il citoplasma sfrutta il TLR4 che si trova a fianco del CD14 così che vengano prodotte citochine proinfiammatorie e pirogeniche, le quali attivano il complemento, la cascata coagulativa e le prostaglandine. Una quantità elevata di LPS può portare a stress respiratori, \underline{coagulazione intravascolare disseminata} e disfunzione multiorgano.\\
Normalmente le endotossine attivano cellule Kupfer, neutrofili, linfociti b e complemento che causano febbre, vasodilatazione, anticorpi e infiammazione.\\
Il \underline{mycobacterium tuberculosis} ha una parete batterica molto complessa con cere e acidi micolici (acidi grassi) che gli conferiscono resistenze immunitarie e antibiotiche, inoltre per questo impiega 12/24h per duplicarsi ed è alcolacido resistente (Zhiel-Neelsen). È per questo che rimane latente per così tanto tempo.
\section*{5. Batteri patogeni}
Principalmente mesofili, eterotrofi e chemiotrofi(?). La virulenza è data da diverse strutture:
\paragraph*{Capsula} Strato saccaridico (proteico per l'antracis) protettivo che permette di resistere alla fagocitosi, a meno che non vi sia opsonizzazione anticorpale o dal complemento, per questo nei vaccini troviamo zuccheri capsulari. Permette anche l'adesione come lo streptococcus mutans (placca). In alcuni batteri (pneumococco) previene il killing intracellulare e non la fagocitosi.\\
È evidenziata dall'inchiostro di china e la sua composizione può portare a mimetismo nel caso di presenza di acido ialuronico (pyogenes) e acido polisialico, una molecola self della Meningitidis che impedisce la funzionalità dei vaccini.
\paragraph*{Pili o fimbrie}Sono strutture proteiche (piline) con terminazioni di adesine che permettono al batterio di ancorarsi ai recettori dell'ospite. Il pilo F è sessuale e viene usato nella coniugazione. La \underline{fragilità dei pili} richiede al batterio di riprodurli frequentemente, l'adesina può essere modificata garantendo maggiore colonizzazione e resistenza immunitaria. \\
Saturando le adesine dei batteri con lo zucchero specifico si può inibire il suo legame alle cellule dell'ospite.
\paragraph*{Biofilm} Un microrganismo planktonico (libera) aderisce ad una superificie inerte o biologica diventanto "sessile", a questo si aggiungono altri microrganismi a formare una colonia molto varia con miceti, batteri, virus. Il biofilm viene protetto da uno strato di \textbf{zuccheri} che impedisce l'attività dei vaccini e protegge la colonia dal sistema immunitario. La colonia può separarsi per colonizzare nuove superfici, questo avviene per rolling, rippling ("ondate"), streaming, seeding (singole cellule), detaching (gruppi di cellule). \\
Sono molto presenti in \textbf{utensili ospedalieri} come cateteri, stent e protesi a causa di Stafilococcus Aureus e Epidermidis, mentre nel caso di fibrosi cistica viene aumentata la crescita dello Pseudomonas Aeroginosae.\\
All'interno del biofilm i batteri comunicano tra loro rilasciando sostanze chimiche, chiamate \underline{autoinduttori}, che inducono agli altri l'espressione di determinati geni. In condizioni normali gli autoinduttori non fanno la differenza, ma nel biofilm la concentrazione cellulare è abbastanza alta, questo meccanismo è il \textbf{Quorum sensing}: regolazione trascrizionale dipendente dalla densità cellulare che può aumentare bioluminescenza, biofilm, spore, pigmenti e altri fattori di virulenza. Alcune sostanze possono inibire e interferire il Quorum sensing.\\
La placca dentale è un biofilm che inizia grazie alla capsula del Mutans, nel caso si dovesse rovinare lo smalto, i batteri raggiungerebbero la dentina causando le carie. Lo smalto è sensibile alle variazioni di pH che avvengono con la produzione di acido lattico da parte dei batteri a causa di un elevato apporto di zuccheri.
\paragraph*{Flagelli} Appendici proteiche (flagellina) che permette al microrganismo di spostarsi, è sempre presente in bacilli, vibrioni e spirilli. Si ancora al batterio grazie a 3 anelli: P(etidoglicano), S(upermembrana) e M(embrana). Nel caso dei gram- è presenta anche l'anello L(PS). Il flagello è l'antigene H.\\
Quando il flagello ruota in senso antiorario procede in avanti, all'opposto si capovolge. Il movimento viene sfruttato basandosi su chemiotassi positiva o negativa da sostanza attraenti o repellenti. 
\subsection*{Tossine} Sostanze tossiche inattivate da formalina o TDDR formando tossoidi. 
\paragraph*{Esotossine} Proteine termolabili e distrutte dai succhi gastrici (tranne le enterotossine stafilococcihe che resistono a entrambi i processi) prodotte sia da gram+ che gram-. \\
Possono essere monomeriche (emolitiche e citolisine che provocano pori nella membrana), \textbf{dimeriche} (porzione tossica A e legante B come tetano e botulino), multimeriche (AB con più sottocomponenti come il colera) e multifattoriali.\\
Hanno trofismo diverso: possono essere neurotrope, enterotossine e pantrope (recettori diffusi).
\end{document}
